\chapter{Ерөнхий хэсэг}
\section{Сэдэв сонгосон үндэслэл}

Гүн сургалтад суурилсан хүний байрлал тодорхойлох технологийн хөгжлөөр энгийн камерын бичлэгээс хүний хөдөлгөөнийг үнэтэй тоног төхөөрөмж, биед зүүдэг мэдрэгчгүйгээр тоон өгөгдөл болгох боломж бүрдсэн. Харин хос бүжиг, спорт, сэргээн засах эмчилгээ зэрэгт олон хүний хөдөлгөөний синхрончлолыг одоог хүртэл мэргэжилтэн нүдээр, субъектив байдлаар үнэлсээр байна. Иймд синхрончлолыг автоматаар, тоон утгаар үнэлэх системийг судалгааны сэдвээ болгон сонгосон.

Хүний байрлал тодорхойлох чиглэлд OpenPose, HRNet, MediaPipe Pose, YOLO-Pose зэрэг олон загвар боловсруулагдсан бөгөөд хөдөлгөөний ижил төстэй байдлыг косинус төстэй байдал, Евклидийн зай, динамик цагийн хазайлт (Dynamic Time Warping, DTW) \cite{sakoe1978}, хөндлөн корреляц зэрэг аргаар харьцуулдаг. Гэвч эдгээр судалгааны ихэнх нь хүнийг лавлах хөдөлгөөнтэй харьцуулдаг бөгөөд зэрэг хөдөлж буй олон хүний синхрончлол, мөн бичлэгийн нягтаршил, кадрын давтамжийн нөлөөг судалсан ажил цөөн байна.

\section{Системийн хэрэгцээ шаардлага}

Систем нь дараах шаардлагуудыг хангах ёстой. Үүнд:

\begin{itemize}
\item
  Хоёр болон түүнээс дээш хүнтэй видеоноос хүн бүрийн гол цэгийг илрүүлж, тухайн хүнд хамааруулах
\item
  Биеийн хэмжээ, камерын зайнаас үл хамааран хөдөлгөөнийг харьцуулах, гол цэгийн шуугианыг бууруулах
\item
  Хос бүрийн болон бүлгийн синхрончлолыг бүхэл бичлэг, хугацааны хэсэг, биеийн хэсгээр тоон утгаар гаргах
\item
  Цагийн хоцрогдлыг тодорхойлж, бүлгээсээ хамгийн их зөрж буй хүнийг илрүүлэх
\item
  Араг яс зурагдсан видеог синхрончлолын үр дүнгийн хамт харуулах
\item
  Видеог вэб хөтөч эсвэл сервер дээр боловсруулж, ижил форматтай үр дүн гаргах
\end{itemize}

\section{Судалгааны аргууд}

Энэхүү судалгаанд дараах аргуудыг ашигласан.

\begin{itemize}
\item
  Онолын судалгаа: HPE загварууд болон хөдөлгөөний ижил төстэй байдлын хэмжүүрүүд
\item
  Харьцуулсан шинжилгээ: HPE загваруудыг гол цэгийн чанар, хурдаар харьцуулах
\item
  Алгоритм боловсруулах: гол цэгийн өгөгдлийг цэвэрлэх, синхрончлолыг тооцоолох
\item
  Туршилт: хяналттай болон бодит бичлэг дээр алгоритмыг шалгах
\item
  Статистик үнэлгээ: дундаж, стандарт хазайлт, дундаж абсолют алдаа (MAE), PCK
\item
  Программ хангамжийн хөгжүүлэлт: системийн архитектур ба хэрэгжилт
\end{itemize}

\section{Ажлын үр дүн, ач холбогдол}

Судалгааны үр дүнд олон хүнтэй видео бичлэгээс хөдөлгөөний синхрончлолын түвшнийг автоматаар үнэлэх систем боловсруулагдсан. Ач холбогдол нь:

\begin{itemize}
\item
  Хөдөлгөөний синхрончлолын субъектив үнэлгээг тоон, давтагдах боломжтой үнэлгээгээр орлуулах боломжийг харуулах
\item
  HPE загварын гүйцэтгэл болон оролтын өгөгдлийн шинж чанарын нөлөөг туршилтаар үнэлэх
\item
  Үнэтэй тоног төхөөрөмж, биед зүүдэг мэдрэгчгүйгээр энгийн камераар хөдөлгөөн шинжлэх хямд шийдлийг харуулах
\end{itemize}
